\section{Experimental evaluation methodology}
\label{sec:jesper-evaluation}
\contributionauthor{Jesper Eggers}

The agents differ in representation, network, mask, rewards, and training
history. Whole-policy comparisons therefore cannot isolate the cause of a
score difference. We distinguish complete-agent performance from controlled
component interventions and engine-based constraint diagnostics.

\subsection{Separate learning reward, selection score, and game performance}
\contributionauthor{Jesper Eggers}
For episode $i$, the primary outcome is the official score
\[
S_i=C_i+5K_i,
\]
where $C_i$ counts collected coins and $K_i$ counts credited opponent kills.
Crate destruction, survival, and threatening an opponent matter strategically,
but do not directly increase this score. We record survival, bomb use,
predicted threats, self-destruction, and episode length as diagnostics, so that a
score change can be interpreted in terms of behavior. For competitive games,
the focal agent's score should also be compared with its opponents' scores.
We define score-leader share by dividing one unit of credit equally among
agents tied for the highest final score. Games continue after the focal
agent dies, so scores include later credits from its bombs and opponents'
subsequent play. All four agents, including those eliminated, enter the
ranking. This measures relative score, not last-survivor success.

The RUEHL checkpoint-selection function in the competitive profile is
\[
J=\overline{S}+5\overline{K}-5\widehat{p}_{\mathrm{self}}
 =\overline{C}+10\overline{K}-5\widehat{p}_{\mathrm{self}}.
\]
Thus, checkpoint selection assigns a kill twice its official point value
relative to a coin. The selection objective $J$ and the final agent's shaped reward both differ
from tournament score and require distinct labels.

Similarly, RUEHL's recorded ``completion'' event requires that the focal
agent has collected the configured total number of coins and that no
opponents remain. In a competitive game, an opponent can collect even one
coin and make this criterion unattainable. A zero completion rate therefore
does not establish a zero survival rate or a zero win rate. This distinction
is necessary when interpreting the curriculum table.

\subsection{Independent tests and controlled comparisons}
\contributionauthor{Jesper Eggers}
The twenty-game greedy evaluations used for checkpoint selection are
validation measurements. Repeatedly choosing the largest score also selects
favorable noise; a monotone best-so-far curve need not show convergence.
Fresh games after freezing the checkpoint estimate subsequent performance.

We evaluated the existing final-agent checkpoint in the classic scenario,
with zero exploration, no learning updates, and 200 fresh episodes per
configuration. The checkpoint was selected by the existing callback's
default filename before observing these results, rather than by test-score selection. Strict tensor verification and
warning-to-error handling exclude silent random fallback. No architecture, reward,
or trained weight was changed.

The four configurations cross the final network's greedy Q-value ranking
or uniform random selection with the final temporal mask or RUEHL mask.
Both learned configurations use the final agent's fourteen spatial and
twenty-two scalar inputs. Thus, ``Final Q + RUEHL mask'' is a deployment
variant of the final network, not the trained RUEHL agent. Each configuration
faces three rule-based opponents and, separately, a mixed group containing
one rule-based, one coin-collector, and one peaceful agent. This yields
1,600 games.

The mask intervention changes the complete implementation, including
escape checks, occupancy handling, and fallback rules; it does not isolate
lookahead alone. The final network also retains its temporal-survival input
features under both masks. Consequently, the mask swap removes neither all
explicit safety reasoning nor all engineered guidance.

All four configurations start each paired episode from the same arena and spawn
assignment, with board seeds separated from checkpoint selection. Merely
launching two multi-round runs with the same seed is insufficient: the
environment also uses its random generator to permute action order, so
different episode lengths consume different numbers of random draws and
alter later boards. We reset the board generator before each round, verify
hashes of the arena, hidden-coin placement, and spawns, then use a separate
seeded generator for action order. Each agent has isolated Python and NumPy
random streams. Episodic position history is cleared between games. The
200 seeds are 270920260--270920459; matching seeds pair configurations
within each opponent group, not their subsequent trajectories.

We report episode-level means and 95\% bootstrap intervals using 10,000
resamples, and paired differences by resampling complete seed-matched
episodes. Survival and self-destruction proportions also receive Wilson
intervals; fractional score-leader share uses the bootstrap.
Twenty episodes
give particularly coarse estimates: two self-destructions correspond to
10\%, but the 95\% Wilson interval is approximately 2.8--30.1\%.
The intervals are pointwise, without correction across the exploratory
component contrasts. Evaluation
uncertainty is a separate issue from variability across independently
trained networks \citep{henderson2018matters,agarwal2021precipice}. With only
one frozen trained artifact in this experiment, conclusions concern that
checkpoint, not a
reliable ranking of learning algorithms.

The mask audit checks candidate actions independently of weights; the
frozen experiment measures selected actions and competitive outcomes.
Neither estimates sample efficiency under a different training mask.
